% Annexe de la partie Immobilier (refonte du 16/11/2026, coupée le 17/11/2026 selon synthese/CONSIGNES_COUPE.md).
% Chiffres : sorties/calculs_foncier.txt et .json (synthese/calculs_foncier.py) ; ../foncier/sorties/simulation.txt
% (foncier/simulation.py) ; registre : ../revue/hypotheses_foncier.md ; figures : ../foncier/figures_note.py. Base 2025
% (décision 5 du 17/11/2026) : données 2025 par ../foncier/donnees_2025.py. Figures n04, n05 et n10 venues du corps.
% Le tableau « base 2024 / base 2025 » de la version du 17/11 (avant coupe) est retiré : chiffres de la base 2024 par
% BASE_FONCIER=2024 (voir calculs_foncier.py) et dans l'historique git (commit e7232a5).

\subsection{Champ, données et calibrage}

\textbf{Champ.} Les 30,5~millions de logements des particuliers de France métropolitaine hors Alsace-Moselle (dont DVF
ne publie pas les ventes), environ 94~\% de la métropole : parc du recensement de 2022 moins les logements HLM loués,
résidences secondaires, logements vacants et environ un million de logements d'autres personnes morales compris. L'unité
de calcul est la cellule « commune × maisons » ou « commune × appartements », dont tous les logements sont supposés
identiques ; la taxe de chaque commune est répartie entre statuts d'occupation au prorata du nombre de logements
(figure~\ref{fig:fon_qui}).

\textbf{Base 2025.} Taxe foncière bâtie par commune (REI 2025), ventes de logements anciens (DVF 2025), comptes de
patrimoine et PIB de 2025 (Insee, comptes de la Nation publiés en 2026), droits et taxe foncière sur les logements (Compte
du logement 2025), droits de toutes les collectivités (OFGL, comptes de gestion 2025). Restent antérieurs le parc
(recensement de 2022), le rapport taxe/valeur de chaque commune (2017) et les prix des logements (ventes de 2023-2024).

\textbf{Valeur.} Prix médian au m² des ventes 2023-2024 de la commune (à défaut, de l'intercommunalité, puis du
département, sous dix ventes) × surface médiane des logements vendus : 6\,762~Md€, 20~\% de moins que les comptes de
patrimoine de l'Insee de fin 2025 ramenés au champ (8\,467~Md€). L'écart n'est pas uniforme : le prix moyen des ventes
dépasse la valeur modélisée de 24~\% dans le dixième de communes le moins cher, de 6 à 10~\% au milieu et de 21~\% dans le
plus cher. Les recettes sont justes, mais les taux en part de la valeur du modèle sont surestimés d'environ un cinquième ;
le corps du texte donne donc les taux en valeurs de l'Insee. Les prix de 2025 sont inférieurs de 1,5~\% à ceux de
2023-2024 (indices Notaires-Insee : Paris $-$2,0~\%, Île-de-France hors Paris $-$3,6~\%, province $-$1,0~\%).

\textbf{Part du sol.} Maisons : prix au m² des terrains à bâtir de la commune (à défaut, celui du département ajusté par
l'écart de prix des maisons, élasticité de 1,54 estimée sur 3\,307~communes) × parcelle médiane plafonnée à 1\,500~m².
Appartements : un moins le rapport entre un coût de reconstruction uniforme (1\,464~€/m², calé pour que le sol fasse
46,2~\% de la valeur, comme dans les comptes de patrimoine des ménages de fin 2025) et le prix au m². Bornes de 10 et
90~\%. Le sol vaut 3\,126~Md€ ; sa part va de 16~\% dans le dixième de logements le moins cher à 68~\% dans le plus cher,
et atteint 85~\% pour l'appartement parisien médian.

\textbf{Taxe foncière actuelle.} Rapport taxe/valeur de la commune en 2017 (André et Meslin, 2023) × valeur × croissance
de la taxe de la commune de 2017 à 2025 (REI), calé sur 19,2~Md€ (2017) × 1,333 (croissance de la taxe foncière bâtie de
la métropole de 2017 à 2025) × part du champ : 24,2~Md€ en 2025, 795~€ par logement, 0,36~\% de la valeur. Les communes
dont la taxe varie de moins de $-$20~\% ou de plus de +30~\% entre 2024 et 2025 (surtout des communes nouvelles,
0,2~\% des logements) gardent leur croissance de 2017 à 2024, prolongée de la croissance nationale de 2025 (+2,8~\%).

\textbf{Droits de mutation.} 5,81~\% (4,50~\% départemental, 1,20~\% communal, frais d'assiette de 2,37~\% de la part
départementale) des 212~Md€ de ventes de logements anciens de 2025 (840\,000 ventes, 2,8~\% du parc) : 12,3~Md€, dont 9,5
départementaux, 2,5 communaux et 0,2 de frais (196~Md€ de ventes et 11,4~Md€ de droits par an en 2023-2024). Rapportés au
parc, 404~€ par logement et par an, 0,18~\% de la valeur. Le relèvement temporaire de la part départementale à 5~\% n'est
pas compté : les droits perçus sur les logements en 2025 (13,6~Md€, France entière ; figure~\ref{fig:fon_cycle})
dépassent de 0,4~Md€ ce qu'aurait donné le taux permanent si le champ en représente la même part qu'en 2023-2024
(93,9~\%), soit un taux moyen effectif d'environ 6,0~\% (calcul de l'auteur, non vérifié).

\textbf{Scénarios.} Tous remplacent la moitié des droits sur les logements (6,2~Md€ sans réaction des ventes) par de la
taxe foncière, à recette constante. Réforme naïve : taxe actuelle relevée de 25,4~\% partout, soit 30,4~Md€ de taxe sur
les logements des particuliers. « N4 » (proposition de la version précédente) : chaque commune garde sa recette, répartie
selon deux taux (sol = 4 × bâti), et un taux national de 0,197~\% sur le sol des logements du modèle remplace les droits.
Proposition : même part communale, droits divisés par deux aussi sur les autres immeubles, taux national sur tous les
terrains privés ; la taxe des logements des particuliers atteint 28,9~Md€.

\figbloc{foncier_n04_qui_paie}{Les propriétaires occupants paient un peu plus de la moitié de la taxe foncière des logements des particuliers}{fig:fon_qui}

\figbloc{foncier_n05_cycle}{Les droits sur les ventes de logements ont perdu un tiers de leur produit entre 2022 et 2024, avant de se redresser en partie en 2025}{fig:fon_cycle}

\subsection{L'assiette du taux national : tous les terrains privés}

Le taux national porte sur le sol sous les bâtiments et les terrains non bâtis des ménages, des sociétés et des
associations, soit 6\,656~Md€ fin 2025, hors terres cultivées (902~Md€) et terrains publics (1\,094~Md€)
(tableau~\ref{tab:fon_insee}) ; le sol sous les bâtiments des ménages en représente 62,1~\%.

\begin{table}[!htbp]\centering\scriptsize
\caption{Valeur des terrains par secteur, fin 2025 (Md€)}\label{tab:fon_insee}
\begin{tabularx}{\linewidth}{@{}>{\raggedright\arraybackslash}X rrrr@{}}
\toprule
Secteur & Sol sous les bâtiments & Terres cultivées & Autres terrains (non bâtis) & Dans l'assiette \\
\midrule
Ménages & 4\,136 & 193 & 452 & 4\,588 \\
Sociétés non financières (dont HLM) & 1\,701 & 710 & 113 & 1\,814 \\
Sociétés financières & 200 & --- & --- & 200 \\
Institutions sans but lucratif & 54 & --- & --- & 54 \\
Administrations publiques & 941 & --- & 153 & exclues \\
Ensemble & 7\,032 & 902 & 717 & 6\,656 \\
\bottomrule
\end{tabularx}
\par\vspace{2pt}{\raggedright\scriptsize Source : Insee, comptes nationaux annuels base 2020, tableau 8.200, millésime
de juin 2026 (\texttt{foncier/donnees/terrains\_secteurs\_2025.csv}), cohérent avec l'Insee Première n° 2135 (novembre
2026 : 4\,190~Md€ de terrains bâtis pour les ménages au sens large). Ce millésime relève de 5~\% la valeur des terrains
bâtis de 2024 (7\,044 au lieu de 6\,729~Md€). Le sol est une valeur résiduelle (valeur de marché moins bâti au coût de
remplacement), sans marge d'incertitude publiée ; les « autres terrains » comprennent des terrains naturels et boisés
valorisés au prix des terres agricoles et des forêts.\par}
\end{table}

\textbf{Méthode.} Le taux national lève $R$ = 7,58~Md€ à l'échelle du champ (6,16~Md€ de droits sur les logements, 1,42
sur les autres immeubles). La part des logements des particuliers, $s = 4\,136/6\,656$, est répartie entre logements au
prorata de la valeur de leur sol dans le modèle : $sR/3\,126 = 0,151~\%$ des valeurs du modèle. En valeurs de l'Insee, le
taux vaut $R/(6\,656 \times 0,946) = 0,120~\%$, si le champ représente la même fraction (94,6~\%) des terrains de chaque
secteur. Les droits sur les autres immeubles (3,13~Md€ en 2025, 3,08~Md€ par an en 2023-2024) sont la différence entre le
total des collectivités (OFGL) et les logements (Compte du logement), deux sources de périmètres différents ; ils sont
ramenés au champ et au taux permanent par le rapport entre droits du modèle et droits perçus sur les logements en 2025
(0,907). Le sol des 4,2~millions de logements HLM loués est estimé à 603~Md€ en valeurs de l'Insee, avec la valeur
médiane des logements de leur commune (un majorant).

\begin{table}[!htbp]\centering\scriptsize
\caption{Le taux national selon l'assiette}\label{tab:fon_sens}
\begin{tabularx}{\linewidth}{@{}>{\raggedright\arraybackslash}X rrrrr@{}}
\toprule
Assiette & Taux (Insee) & Taux (modèle) & Paris & Dernier dixième & Hausses $>$ 300~€ \\
\midrule
Sol des logements seul, droits des logements seuls (N4) & 0,157~\% & 0,197~\% & +129~\% & +602~€ & 25,1~\% \\
Terrains bâtis privés & 0,131~\% & 0,165~\% & +108~\% & +500~€ & 20,7~\% \\
Tous les terrains privés (proposition) & 0,120~\% & 0,151~\% & +99~\% & +456~€ & 19,0~\% \\
Tous les terrains privés, sol des HLM exonéré & 0,132~\% & 0,166~\% & +109~\% & +503~€ & 20,9~\% \\
Tous les terrains privés, terres agricoles comprises & 0,106~\% & 0,133~\% & +87~\% & +399~€ & 15,8~\% \\
Tous les terrains privés, droits des seuls logements & 0,098~\% & 0,122~\% & +80~\% & +366~€ & 14,3~\% \\
\bottomrule
\end{tabularx}
\par\vspace{2pt}{\raggedright\scriptsize Sans réaction des ventes ; avec une reprise de 6~\% des ventes par point, taux
et hausses baissent d'environ 18~\%. Paris : variation moyenne de la taxe des logements ; dernier dixième : variation
moyenne par logement dans le dixième le plus cher. Les terres agricoles toucheraient des exploitants déjà soumis à la
taxe sur le foncier non bâti.\par}
\end{table}

\subsection{Résultats complets par dixième de logements}

\begin{table}[!htbp]\centering\scriptsize
\caption{Taxe foncière et droits par dixième de logements, selon le prix au m² de leur commune}\label{tab:fon_dec}
\setlength{\tabcolsep}{3.5pt}
\begin{tabularx}{\linewidth}{@{}c r *{4}{r} *{3}{r} *{3}{r}@{}}
\toprule
& Prix & \multicolumn{4}{c}{Taxe / valeur (\%)} & \multicolumn{3}{c}{Variation de taxe (€/an)} & \multicolumn{3}{c}{Taxe moins droits (€/an)} \\
\cmidrule(lr){3-6}\cmidrule(lr){7-9}\cmidrule(l){10-12}
Dixième & (€/m²) & auj. & naïve & N4 & prop. & naïve & N4 & proposition & naïve & N4 & prop. \\
\midrule
1 & 1\,035 & 0,54 & 0,67 & 0,55 & 0,55 & +114 & +16 & +10 (+2~\%) & +33 & $-$65 & $-$72 \\
2 & 1\,404 & 0,49 & 0,61 & 0,52 & 0,51 & +147 & +36 & +25 (+4~\%) & +42 & $-$69 & $-$80 \\
3 & 1\,693 & 0,49 & 0,61 & 0,53 & 0,52 & +175 & +53 & +38 (+5~\%) & +45 & $-$77 & $-$92 \\
4 & 2\,000 & 0,46 & 0,58 & 0,53 & 0,51 & +193 & +104 & +81 (+11~\%) & +47 & $-$42 & $-$65 \\
5 & 2\,356 & 0,43 & 0,54 & 0,51 & 0,49 & +212 & +144 & +113 (+14~\%) & +47 & $-$21 & $-$52 \\
6 & 2\,737 & 0,41 & 0,52 & 0,50 & 0,48 & +224 & +177 & +137 (+16~\%) & +36 & $-$10 & $-$50 \\
7 & 3\,176 & 0,39 & 0,49 & 0,49 & 0,47 & +235 & +236 & +185 (+20~\%) & +26 & +26 & $-$25 \\
8 & 3\,749 & 0,36 & 0,45 & 0,46 & 0,44 & +252 & +295 & +228 (+23~\%) & +20 & +63 & $-$4 \\
9 & 4\,746 & 0,31 & 0,38 & 0,41 & 0,39 & +256 & +358 & +274 (+27~\%) & $-$38 & +63 & $-$21 \\
10 & 6\,261 & 0,18 & 0,23 & 0,31 & 0,28 & +212 & +602 & +456 (+55~\%) & $-$258 & +132 & $-$14 \\
\bottomrule
\end{tabularx}
\par\vspace{2pt}{\raggedright\scriptsize Dixièmes de logements (et non de communes), classés selon le prix médian au m²
de leur commune. Taxe moins droits : baisse des droits rapportée à l'ensemble des logements de la commune. Sans réaction
des ventes.\par}
\end{table}

\textbf{Distribution individuelle (proposition).} Variation de la taxe foncière par logement : $-$9~€ au
10\textsuperscript{e} centile, +36~€ au 25\textsuperscript{e}, +110~€ à la médiane, +238~€ au 75\textsuperscript{e},
+422~€ au 90\textsuperscript{e}, +545~€ au 95\textsuperscript{e} et +647~€ au 99\textsuperscript{e} ; 7,2~\% des
logements dépassent +500~€, 0,9~\% +660~€ et 0,2~\% +1\,000~€ ; 11~\% voient leur taxe baisser (un sur neuf). Avec le taux
national sur le seul sol des logements, les mêmes centiles valent +4, +51, +141, +301, +529, +709 et +800~€
(figure~\ref{fig:fon_distrib}). Plus fortes hausses (cellules d'au moins 1\,000~logements) : maisons de Saint-Tropez
+2\,576~€ (+71~\%), de Megève +1\,927~€, appartements de Saint-Jean-Cap-Ferrat +1\,887~€, maisons de Ramatuelle +1\,799~€,
de La Clusaz +1\,742~€ et de Cassis +1\,673~€ (+109~\%). Hors Paris, le dixième le plus cher prend 40~\%. Par type, les
maisons prennent +112~€ et les appartements +212~€ en moyenne, mais +197 et +97~€ si l'on mesure le sol des maisons par
résidu, comme celui des appartements (coût de reconstruction recalé à 1\,973~€/m²).

\textbf{Vieillissement de la taxe actuelle.} Appliquer le rapport taxe/valeur de 2017 aux valeurs de 2023-2024, puis la
croissance de la taxe de 2017 à 2025, compte deux fois la hausse des prix de 2017 à 2023-2024. Une correction à trois
zones (indices Notaires-Insee : Paris ×1,113, Île-de-France hors Paris ×1,134, province ×1,281) ramène la hausse
parisienne de +99 à +88~\% et celle du dernier dixième de +55 à +52~\% ; la part des hausses de plus de 300~€ ne change
pas (19,0~\%). Une correction propre demande les valeurs de 2017 par commune.

\figbloc{foncier_n12_variantes}{Toutes les variantes de répartition : les deux taux nationaux inversent la régressivité, la taxe relevée l'aggrave}{fig:fon_var_toutes}

\begin{table}[!htbp]\centering\scriptsize
\caption{Toutes les variantes de répartition (droits sur les logements divisés par deux)}\label{tab:fon_variantes}
\begin{tabularx}{\linewidth}{@{}>{\raggedright\arraybackslash}X rrrr@{}}
\toprule
Variante & Premier dixième & Dernier dixième & Hausses $>$ 300~€ & Taxes qui doublent \\
\midrule
Taxe actuelle relevée (réforme naïve) & +114~€ (+25~\%) & +212~€ (+25~\%) & 14~\% & 0~\% \\
Valeur de marché, taux national unique & $-$72~€ ($-$16~\%) & +1\,247~€ (+149~\%) & 28~\% & 11,8~\% \\
Deux taux nationaux, sol = 2 × bâti & $-$149~€ ($-$33~\%) & +1\,559~€ (+187~\%) & 28~\% & 12,5~\% \\
Deux taux nationaux, sol = 4 × bâti & $-$213~€ ($-$47~\%) & +1\,821~€ (+218~\%) & 29~\% & 13,7~\% \\
Sol seul, taux national & $-$315~€ ($-$70~\%) & +2\,236~€ (+268~\%) & 29~\% & 16,3~\% \\
Deux taux par commune (2 × bâti), taux départemental sur le sol & +51~€ (+11~\%) & +463~€ (+55~\%) & 23~\% & 1,2~\% \\
Deux taux par commune (4 × bâti), taux départemental sur le sol & +45~€ (+10~\%) & +454~€ (+54~\%) & 25~\% & 1,9~\% \\
Deux taux par commune, taux national sur le sol des logements (N4) & +16~€ (+4~\%) & +602~€ (+72~\%) & 25~\% & 5,9~\% \\
N4 avec reprise des ventes (+6~\% par point) & +11~€ (+2~\%) & +487~€ (+58~\%) & 20~\% & 4,8~\% \\
Proposition (tous les terrains, tous les immeubles) & +10~€ (+2~\%) & +456~€ (+55~\%) & 19~\% & 4,5~\% \\
\bottomrule
\end{tabularx}
\par\vspace{2pt}{\raggedright\scriptsize Variation moyenne de la taxe foncière par logement. Taux nationaux, en part de
la valeur du modèle : 0,614 et 0,307~\% (sol = 2 × bâti), 0,753 et 0,188~\% (4 ×), 0,972~\% (sol seul). Taux
départemental : de 0,106~\% (Paris : 0,130~\%) à 0,789~\% (Creuse). Sources : \texttt{foncier/sorties/simulation.txt},
\texttt{sorties/calculs\_foncier.txt}.\par}
\end{table}

\figbloc{foncier_n13_distribution}{Le taux national sur tous les terrains réduit les plus fortes hausses d'un cinquième par rapport au taux sur le seul sol des logements}{fig:fon_distrib}

\subsection{Perte sèche}

Deux triangles (méthode : \texttt{revue/methode\_perte\_seche.md}). Sur les ventes, $\sigma V(\tau_1^2-\tau_0^2)/2$, où
$V$ = 212~Md€ de ventes (2025) et $\sigma$ la baisse des ventes par point de droits. Sur le capital bâti,
$\varepsilon_K (K/u)(t_1^2-t_0^2)/2$, où $t$ est l'impôt sur le bâti rapporté à sa valeur ($K$ = 3\,636~Md€), $u$ le
coût d'usage (4,5~\%) et $\varepsilon_K$ l'élasticité du capital logement (1). Le sol n'entre dans aucun triangle. Ces
pertes sèches sont statiques ; à long terme, que les ventes perdues en 2014 ne semblent pas l'avoir été durablement
(Bérard et Trannoy, 2018) plaide pour le bas de la fourchette.

\begin{table}[!htbp]\centering\scriptsize
\caption{Perte sèche évitée (Md€ par an ; signe moins = gain)}\label{tab:fon_ps}
\begin{tabularx}{\linewidth}{@{}>{\raggedright\arraybackslash}X r@{}}
\toprule
Ventes de logements, droits divisés par deux : $\sigma$ = 0 / 3 / 6 / 8 / 15 & 0 / $-$0,80 / $-$1,61 / $-$2,14 / $-$4,02 \\
Ventes de logements, baisse de 30~\% / de 100~\% ($\sigma$ = 6) & $-$1,09 / $-$2,14 \\
Bâti, proposition (impôt sur le bâti de 14,3 à 7,3~Md€) : $\varepsilon_K$ = 1 / 0,5 ; $u$ = 3,5 / 6~\% & $-$0,46 / $-$0,23 ; $-$0,59 / $-$0,34 \\
Bâti, rapport sol/bâti de 2 / 4 / sol seul (impôt sur le bâti de 10,7 / 7,3 / 0~Md€) & $-$0,27 / $-$0,46 / $-$0,62 \\
Bâti, réforme naïve (impôt sur le bâti de 14,3 à 17,9~Md€) & +0,36 \\
Bâti, en comptant les droits comme un impôt annuel sur la détention (0,18~\%) : proposition / réforme naïve & $-$0,98 / +0,04 \\
Ventes des autres immeubles ($\sigma$ = 6 transposé, 49~Md€ de ventes ; ordre de grandeur) & $-$0,37 \\
\midrule
Total, réforme naïve & $-$1,25 \\
Total, proposition, logements / avec les autres immeubles & $-$2,06 / $-$2,44 \\
Variante : coût des logements mal adaptés (Määttänen et Terviö, 2022) à la place du triangle sur les ventes & environ $-$3,8 \\
\bottomrule
\end{tabularx}
\par\vspace{2pt}{\raggedright\scriptsize La variante de Määttänen et Terviö applique aux 30,5~millions de logements le
gain par ménage de leur modèle d'Helsinki (euros de 2004), interpolé entre 30~€ par an pour un droit de 2~\% et 180~€ pour
6~\% (109~€ par logement, 3,3~Md€) ; elle remplace le triangle sur les ventes, qui mesure les mêmes échanges manqués.
Le loyer que se verse un propriétaire occupant n'étant pas imposé, alléger l'impôt sur le bâti peut aussi accroître une
distorsion existante : le signe du terme bâti n'est pas certain.\par}
\end{table}

\subsection{Capitalisation, report de paiement et transition}

\textbf{Capitalisation.} Un impôt annuel permanent $\Delta$ réduit la valeur d'un bien de $\Delta/r$, avec $r$ le taux
d'actualisation réel (3~\%, mesuré à 3-4~\% en Angleterre par Koster et Pinchbeck, 2022). Cas symétrique : la valeur
varie de $-(\Delta T + \Delta D)/r$, où $\Delta D$ est la baisse annuelle moyenne des droits par logement de la commune ;
cas asymétrique : de $-\Delta T/r$. Une montée linéaire sur $n$ années laisse supporter à l'annonce une part
$\sum_t \min(t/n, 1)(1+r)^{-t} \big/ \sum_t (1+r)^{-t}$ de la perte : 97~\% sur trois ans, 94~\% sur cinq, 88~\% sur dix,
77~\% sur vingt.

\begin{table}[!htbp]\centering\scriptsize
\caption{Variation de la valeur des logements des particuliers selon l'incidence (r = 3~\%)}\label{tab:fon_cap}
\begin{tabularx}{\linewidth}{@{}>{\raggedright\arraybackslash}p{3.4cm}*{4}{>{\raggedright\arraybackslash}X}@{}}
\toprule
& Ensemble & Somme des pertes & Appartement parisien médian & Logements perdant plus de 5~\% \\
\midrule
Proposition, cas symétrique & +48~Md€ (+0,7~\%) & 35~Md€ & $-$2\,970~€ ($-$0,7~\%) & 0,9~\% \\
Proposition, cas asymétrique & $-$157~Md€ ($-$2,3~\%) & 168~Md€ & $-$18\,152~€ ($-$4,4~\%) & 5,9~\% \\
Réforme naïve, cas asymétrique & $-$205~Md€ ($-$3,0~\%) & 205~Md€ & $-$4\,466~€ ($-$1,1~\%) & --- \\
\bottomrule
\end{tabularx}
\par\vspace{2pt}{\raggedright\scriptsize Cas symétrique : la hausse de taxe et la baisse des droits passent toutes deux
dans les prix ; cas asymétrique : seule la hausse de taxe y passe.\par}
\end{table}

\textbf{Report de paiement.} La hausse $\Delta$ reportée chaque année pendant $h$ ans au taux $i$ forme une dette de
$\Delta[(1+i)^h - 1]/i$. Pour l'appartement parisien médian ($\Delta$ = 545~€) : 6\,200~€ après dix ans, 10\,100~€ après
quinze, 14\,600~€ après vingt à 3~\% (1,5, 2,4 et 3,5~\% de la valeur) ; 9\,400~€ à 2~\% et 10\,900~€ à 4~\% après quinze
ans. Le taux de 3~\% est une convention : la règle est de facturer le taux d'emprunt de l'État.

\begin{table}[!htbp]\centering\scriptsize
\caption{Calendrier de la proposition sur cinq ans}\label{tab:fon_cal}
\begin{tabularx}{\linewidth}{@{}>{\raggedright\arraybackslash}X rrrrr@{}}
\toprule
Année & 1 & 2 & 3 & 4 & 5 \\
\midrule
Taux national (\% de la valeur du sol, Insee) & 0,024 & 0,048 & 0,072 & 0,096 & 0,120 \\
Recette du taux national (Md€) & 1,5 & 3,0 & 4,5 & 6,1 & 7,6 \\
Droits non remplacés (Md€) & 6,1 & 4,5 & 3,0 & 1,5 & 0 \\
Taxe de l'appartement parisien médian (€, 527 en 2025) & 636 & 745 & 854 & 963 & 1\,072 \\
\bottomrule
\end{tabularx}
\par\vspace{2pt}{\raggedright\scriptsize Droits divisés par deux dès l'annonce (année 0, dont les quelques mois avant le
1\textsuperscript{er}~janvier ne sont pas comptés) ; ventes et taxe foncière de 2025, en euros de 2025. Déficit cumulé :
15,2~Md€ (12,4 avec la reprise des ventes) ; sur trois ans 7,6 (6,2), sur dix ans 34,1 (27,8). Si la montée ne commence
qu'après trois années de droits baissés : 30,3 / 37,9 / 56,8~Md€ (montée sur 3, 5 ou 10 ans) ; avec une baisse des droits
en deux temps (un quart à l'annonce, un quart au démarrage), qui limite ce risque mais recrée l'attente : 18,9 / 26,5 /
45,5~Md€.\par}
\end{table}

\textbf{Autres calendriers.} Avec le taux national plein dès la première année, l'appartement parisien médian prendrait
d'un coup +534~€ (+698~€ avec le taux sur le seul sol des logements). Financer l'attente par une hausse uniforme de la
taxe actuelle demanderait 31~\% de hausse et porterait la taxe à 0,70~\% de la valeur dans le dixième le moins cher,
contre 0,24~\% dans le plus cher.

\textbf{Collectivités.} Droits perdus sur les logements (ventes de 2025, taux permanent) : 6,16~Md€, dont 4,77 pour les
départements, 1,27 pour les communes et 0,11 de frais de l'État ; s'y ajoutent 1,42~Md€ sur les autres immeubles. Reversé
à chaque département à hauteur de ses pertes, le taux national crée une péréquation : sur la seule part des logements,
Paris paie 621~M€ et perd 537~M€ (+83~M€), les Hauts-de-Seine +45~M€, la Creuse $-$4~M€ ; sept départements sont
contributeurs nets, pour 204~M€. Avec le taux sur le seul sol des logements, Paris paierait 811~M€ (+274~M€), les
Hauts-de-Seine +137~M€, et les transferts atteindraient 869~M€.

\subsection{Cas-types complets}

\begin{table}[!htbp]\centering\scriptsize
\caption{Le jeune ménage qui achète le logement médian de sa commune}\label{tab:fon_primo}
\setlength{\tabcolsep}{4pt}
\begin{tabularx}{\linewidth}{@{}>{\raggedright\arraybackslash}X rrrrrr@{}}
\toprule
Logement & Sol & Droits en moins & Taxe (€/an) & dont national & Revend après 7 / 15 ans & Économie rattrapée \\
\midrule
Paris, appartement & 85~\% & 12\,083 & +545 & 534 & +8\,690 / +5\,582 & 37 ans (naïve : jamais) \\
Lyon, appartement & 69~\% & 8\,105 & +303 & 288 & +6\,217 / +4\,488 & 55 ans (jamais) \\
Bordeaux, appartement & 66~\% & 6\,469 & +344 & 223 & +4\,326 / +2\,362 & 28 ans (jamais) \\
Nantes, maison & 46~\% & 13\,051 & +143 & 314 & +12\,161 / +11\,346 & jamais (65 ans) \\
Limoges, maison & 23~\% & 4\,988 & +253 & 61 & +3\,410 / +1\,965 & 30 ans (22 ans) \\
Saint-Étienne, appartement & 10~\% & 2\,266 & $-$42 & 12 & +2\,527 / +2\,766 & taxe en baisse (18 ans) \\
Annecy, maison & 30~\% & 20\,125 & $-$204 & 314 & +21\,394 / +22\,558 & taxe en baisse (jamais) \\
Arcachon, maison & 42~\% & 22\,700 & +123 & 496 & +21\,934 / +21\,233 & jamais (jamais) \\
Guéret, maison & 11~\% & 3\,193 & +27 & 18 & +3\,028 / +2\,876 & jamais (16 ans) \\
Ussel, maison & 10~\% & 2\,982 & +15 & 15 & +2\,886 / +2\,798 & jamais (13 ans) \\
\bottomrule
\end{tabularx}
\par\vspace{2pt}{\raggedright\scriptsize En euros ; taxe de 2025, droits au taux permanent, logement médian aux prix de
2023-2024. Revend : économie de droits moins la valeur actuelle (3~\%) de la hausse de taxe payée pendant la détention,
sans effet sur les prix. Économie rattrapée : durée au bout de laquelle la hausse actualisée égale l'économie (« jamais »
si la hausse perpétuelle actualisée reste inférieure) ; entre parenthèses, la réforme naïve. Part communale : variation
totale moins part nationale. Une famille nantaise qui passe d'un appartement à une maison de 450\,000~€ paie 13\,050~€ de
droits au lieu de 26\,100~€.\par}
\end{table}

\figbloc{foncier_n10_primo}{Un ménage qui achète le logement médian gagne 2\,000 à 13\,000~€ de droits ; la hausse de taxe foncière en reprend moins de la moitié en dix ans}{fig:fon_primo}

\subsection{Changements par rapport à la version précédente}

La version précédente reposait sur la taxe foncière de 2024, les ventes de 2023-2024 et les comptes de patrimoine de fin
2024. En base 2025, les ventes de logements anciens dépassent de 8~\% la moyenne de 2023-2024, la taxe foncière bâtie de
2,8~\% celle de 2024, et les terrains privés valent 4,9~\% de plus, surtout parce que l'Insee a relevé d'environ 5~\% ses
valeurs de 2024. Le taux national reste à 0,12~\%, la répartition des hausses ne change presque pas, et la bascule « taxe
moins droits » devient favorable dans tous les dixièmes ; les conclusions ne changent pas. Quatre chiffres de la version
précédente étaient faux (registre \texttt{revue/hypotheses\_foncier.md}, §~5.2) : la taxe foncière sur les logements
de 2024 valait 30,5~Md€, et non 31,6 (chiffre de 2025) ; le sol seul évite 0,6~Md€ de perte sèche sur le bâti, et non
0,5 ; la péréquation parisienne (« 213~M€ ») supposait un taux national limité à la part départementale des droits ;
et les « 660~€ » attribués à l'appartement parisien médian étaient le seuil des 5~\% de plus fortes hausses. Les dixièmes
portent désormais sur le prix au m² des communes et non sur leur richesse ; le coût des logements mal adaptés remplace le
triangle sur les ventes au lieu de s'y ajouter ; le taux national, qui porte sur tous les terrains privés, monte
progressivement, et la transition n'est plus financée par la taxe assise sur les valeurs de 1970.

\subsection{Hypothèses fortes et chiffres non vérifiés}

\begin{table}[!htbp]\centering\scriptsize
\caption{Hypothèses fortes}\label{tab:fon_hyp}
\begin{tabularx}{\linewidth}{@{}>{\raggedright\arraybackslash}p{4.3cm}>{\raggedright\arraybackslash}p{2.6cm}>{\raggedright\arraybackslash}X@{}}
\toprule
Hypothèse & Où elle joue & Ce qu'elle change si elle est fausse \\
\midrule
Dans une commune, taxe actuelle proportionnelle à la valeur de marché & variations individuelles & Gagnants et perdants plus nombreux et plus marqués (écarts entre quartiers : voir le corps) ; non chiffrable sans les fichiers fonciers \\
Logements d'une cellule commune × type identiques & parts de logements, cas-types & Parts de cellules : ordres de grandeur, pas des nombres d'avis \\
Valeur = prix médian × surface médiane (20~\% sous l'Insee, biais en U) & taux en part de la valeur & Taux du modèle surestimés d'un cinquième ; recettes justes ; régressivité confirmée après correction (environ 0,43 à 0,15~\%) \\
Sol des appartements : coût de reconstruction uniforme (1\,464~€/m²) & part nationale, bâti & 9~\% des logements à la borne de 10~\% ; avec 2\,500~€/m² à Paris, hausse parisienne inférieure d'environ 11~\% \\
Sol des maisons = terrains à bâtir × parcelle médiane & maisons / appartements & Mesuré par résidu, le partage s'inverse : ne rien en conclure \\
Taxe actuelle vieillie en comptant deux fois la hausse des prix & taxe actuelle par commune & Paris +88~\% au lieu de +99~\% ; dixièmes peu touchés \\
Calage sur 19,2~Md€ (2017), champ plus étroit que le modèle & hausses en pourcentage & Taxe actuelle peut-être sous-estimée de 7~\% au plus : hausses en pourcentage un peu surestimées \\
Ventes de 2025 au taux permanent de 5,81~\% & coût, taux national & Au niveau de 2021-2022, la baisse sur les logements coûterait près de 8~Md€ ; au taux temporaire de 6,32~\%, 9~\% de plus \\
Ventes en hausse de 6~\% par point de baisse (convention) & gain, coût & Sans réaction durable, 0,5~Md€ de gain (bâti) seulement ; gain linéaire en $\sigma$ \\
Incidence : trois cas polaires & capitalisation & Études contradictoires (acheteur en France en 2014, vendeur en Allemagne ; capitalisation complète ou nulle au Danemark) \\
Bâti : $\varepsilon_K$ = 1, $u$ = 4,5~\%, construction non bloquée par le droit des sols & gain sur le bâti & De 0,2 à 0,6~Md€ ; plus faible là où le zonage bloque la construction \\
Valeurs de l'Insee des terrains par secteur, même part de champ par secteur & assiette, taux national & Valeurs résiduelles, révisées de 9~\% (2023) puis de 5~\% (2024) ; terrains d'entreprise surestimés : taux et charge des logements plus élevés (sol bâti seul : 0,13~\%) \\
Charge sur les entreprises et les terrains nus non localisée & péréquation & Géographie des 2,9~Md€ non résidentiels inconnue : péréquation calculée pour les seuls logements \\
Droits sur les autres immeubles = différence de deux sources ; même réaction des ventes & coût (1,42~Md€), gain (0,4~Md€) & Montant et gain incertains ; variante limitée aux logements : 0,098~\% \\
Chaque commune garde sa recette ; locaux professionnels, HLM et taxe d'ordures ménagères hors modèle & part communale & Étendre la base aux locaux professionnels créerait des transferts non simulés \\
Valeurs provisoires du sol prêtes dès la première année & calendrier, déficit & Trois ans de retard : +23~Md€ de déficit ; validité juridique non expertisée \\
Actualisation et intérêt du report à 3~\% réels & valeurs actuelles, dette reportée & À 4~\%, la perte actualisée baisse d'un quart \\
Exonérations actuelles (plus de 75~ans modestes, neuf) non modélisées & retraités, taux national & À décider pour le taux national \\
Millésimes différents (taxe/valeur 2017, recensement 2022, prix 2023-2024, taxe et ventes 2025) & tout & Prix de 2025 inférieurs de 1,5~\% ; à actualiser pour une réforme en 2028-2030 \\
\bottomrule
\end{tabularx}
\end{table}

\begin{table}[!htbp]\centering\scriptsize
\caption{Chiffres non vérifiés sur une source primaire}\label{tab:fon_nv}
\begin{tabularx}{\linewidth}{@{}>{\raggedright\arraybackslash}p{5.2cm}>{\raggedright\arraybackslash}p{3.2cm}>{\raggedright\arraybackslash}X@{}}
\toprule
Chiffre & Où il sert & État \\
\midrule
Part immobilière des successions : 60~\% (12,7~Md€ en 2025) & inventaire & clé du CPO, aucune ventilation officielle \\
Taxe foncière hors logements : 11,5~Md€ en 2025 & inventaire & REI moins Compte du logement \\
Droits sur les autres biens : 3,1~Md€ en 2025 & coût de la baisse & OFGL moins Compte du logement \\
Frais d'assiette de l'État : 2,37~\% de la part départementale & taux de 5,81~\% & CPO et BOFiP ; article du CGI non relu \\
Taxe d'ordures ménagères assise sur la base de la taxe foncière & transition & article 1522 du CGI non relu \\
Plafond de revenu de l'exonération des plus de 75~ans (12\,793~€ pour une part en 2026) & retraités & source secondaire \\
Quotité moyenne financée de 84,7~\% (apport de 15~\%) & poids des droits dans l'apport & CCSF 2026 d'après la note de lecture \\
Taxe foncière non récupérable sur le locataire, déductible au régime réel & cas de la bailleuse & décret n°~87-713 et article 31 du CGI non relus \\
Coût de reconstruction de 1\,464~€/m² & part du sol & calé, jamais comparé à une source publiée \\
Produit du relèvement temporaire des droits en 2025 : 0,4~Md€ & droits de 2025 & suppose que le champ pèse la même part des droits qu'en 2023-2024 \\
Plus-values immobilières (2,8~Md€), contribution de sécurité immobilière (0,7~Md€), taxe sur les bureaux d'Île-de-France (1,1~Md€) & inventaire & chiffres de 2024 ; plus-values d'après l'audition de la Cour des comptes (rapport non lu) \\
Sol des HLM : 603~Md€ & variante d'exonération & estimation du modèle, majorant \\
Réaction des ventes : 6~\% par point & gain, coût & convention construite sur la littérature \\
Danemark : base de 80~\% de la valeur depuis 2024, report automatique des hausses & report par défaut & textes officiels lus via des outils de recherche (\texttt{lvt/notes/transitions\_et\_saillance.md}) \\
Territoire de la capitale australienne : transition ralentie & calendrier & presse et organismes favorables à la réforme \\
Recours constitutionnels pendants en Allemagne & expériences étrangères & source associative \\
Karlsruhe : maisons +76~\%, copropriétés $-$20~\% & expériences étrangères & calcul de la ville (présentation du 7/10/2024) \\
Banzhaf et Lavery : +3~points de croissance décennale du nombre de pièces & problème du bâti & version de travail de 2008 \\
Recours au report facultatif inférieur à 1~\% aux États-Unis & report par défaut & Baer (1998), cité par Shan (2010) \\
\bottomrule
\end{tabularx}
\end{table}
